% Adding one field, in a shared graph and in a federated one.
%
% This is the book's first figure and therefore fixes the idiom SPEC decision
% 33 describes: square corners, no fills, one stroke weight, sans-serif labels,
% black on white, and order shown as a timeline rather than as a box-and-arrow
% architecture diagram. Later figures are matched against this file, not
% against the decision's wording.
%
% Bare tikzpicture (house style): the figure environment, caption and label
% live at the call site in the section file.
\begin{tikzpicture}[
  x=1mm, y=1mm,
  every node/.append style={font=\sffamily\scriptsize},
  axis/.style={draw, line width=0.4pt, -{Straight Barb[length=1.4mm, width=1.6mm]}},
  tick/.style={draw, line width=0.4pt},
  event/.style={anchor=south, align=center, text width=33mm, inner sep=1pt},
  span/.style={draw, line width=0.4pt},
  spanlabel/.style={anchor=north, align=center, inner sep=2pt},
  lane/.style={anchor=west, font=\sffamily\scriptsize\itshape},
]

% ------------------------------------------------------------------- lane A
\node[lane] at (0,13) {A. one schema, owned by one team};

\draw[axis] (0,0) -- (150,0);
\foreach \x in {15,55,95,135}{\draw[tick] (\x,-1.5) -- (\x,1.5);}

\node[event] at (15,3)  {you ask for\\the field};
\node[event] at (55,3)  {it joins\\their queue};
\node[event] at (95,3)  {they edit the schema\\and the resolver};
\node[event] at (135,3) {one deployment,\\everyone's};

\draw[span] (48,-4) -- (48,-7) -- (102,-7) -- (102,-4);
\node[spanlabel] at (75,-8) {another team's calendar, and you are not on it};

% ------------------------------------------------------------------- lane B
\begin{scope}[shift={(0,-42)}]
  \node[lane] at (0,13) {B. a federated graph};

  \draw[axis] (0,0) -- (150,0);
  \foreach \x in {20,75,130}{\draw[tick] (\x,-1.5) -- (\x,1.5);}

  \node[event] at (20,3)  {you edit your\\own subgraph};
  \node[event] at (75,3)  {it composes,\\or it does not};
  \node[event] at (130,3) {your deployment,\\yours alone};

  \draw[span] (64,-4) -- (64,-7) -- (86,-7) -- (86,-4);
  \node[spanlabel] at (75,-8) {a machine, and it answers every time you ask};
\end{scope}

\end{tikzpicture}
